\documentclass[11pt]{article}

\usepackage[a4paper,margin=2.5cm]{geometry}
\usepackage[utf8]{inputenc}
\usepackage[T2A, T1]{fontenc}
\usepackage[ukrainian,english]{babel}
\usepackage{amsmath,amssymb}
\usepackage{graphicx}
\usepackage{hyperref}
\usepackage{listings}
\usepackage{xcolor}

% Monospace font for code, type names, functions, etc.
\usepackage{inconsolata}
\newcommand{\code}[1]{\texttt{#1}}

\usepackage{newunicodechar}
\newunicodechar{ʼ}{'}

% Big-O and big-Theta notation
\newcommand{\BigO}[1]{\mathcal{O}(#1)}
\newcommand{\BigTheta}[1]{\Theta(#1)}

\lstset{
    basicstyle=\ttfamily\small,
    breaklines=true,
    frame=single,
    columns=fullflexible,
    backgroundcolor=\color{gray!10},
    keywordstyle=\color{blue},
    commentstyle=\color{gray},
    stringstyle=\color{teal},
}

\lstdefinelanguage{Rust}{
    keywords={as,break,const,continue,else,enum,fn,for,if,impl,in,let,loop,match,mod,mut,pub,return,self,Self,struct,use,where,while},
    sensitive=true,
    comment=[l]{//},
    string=[b]",
}

\lstdefinestyle{ruststyle}{
    language=Rust,
    basicstyle=\ttfamily\footnotesize,
    numbers=left,
    numberstyle=\tiny\color{gray},
    numbersep=8pt,
    tabsize=4,
    showstringspaces=false,
    frame=single,
    breaklines=true,
    columns=fullflexible,
    backgroundcolor=\color{gray!10},
    keywordstyle=\color{blue!70!black}\bfseries,
    commentstyle=\color{gray},
    stringstyle=\color{teal},
}

\title{Спеціальні розділи обчислювальної математики\\
Компʼютерний практикум №2\\[1em]
Багаторозрядна модулярна арифметика}
\author{ФІ-43 Редько Микита}
\date{\today}

\begin{document}
\selectlanguage{ukrainian}

\maketitle

\section{Мета роботи}
Отримання практичних навичок програмної реалізації багаторозрядної арифметики;
ознайомлення з прийомами ефективної реалізації критичних по часу
ділянок програмного коду та методами оцінки їх ефективності.

\section{Реалізація}

Бібліотеку для роботи з $m$-бітними цілими числами ($m = 2048$), створену
на компʼютерному практикумі №1, доповнено такими операціями:
\begin{itemize}
    \item \code{gcd\_big\_int}, \code{lcm\_big\_int} "--- обчислення НСД
          (алгоритм Евкліда) та НСК ($\frac{a}{\gcd(a, b)} \cdot b$) двох чисел;
    \item \code{barrett\_precompute\_big\_int} "--- одноразове обчислення
          $\mu = \left\lfloor \beta^{2k} / n \right\rfloor$ для модуля $n$,
          де $\beta = 2^{64}$, а $k$ "--- кількість розрядів $n$;
    \item \code{barrett\_reduction\_big\_int} "--- редукція Барретта $x \bmod n$
          для $x < \beta^{2k}$ (якщо $x < n$, число повертається без обчислень);
    \item \code{add\_mod\_big\_int}, \code{subtract\_mod\_big\_int} "--- додавання
          та віднімання чисел за модулем (з одним умовним відніманням або додаванням $n$);
    \item \code{mult\_mod\_big\_int}, \code{square\_mod\_big\_int} "--- множення
          чисел та піднесення чисел до квадрата за модулем;
    \item \code{horner\_exponentiation\_mod\_big\_int} "--- піднесення числа
          до багаторозрядного степеня $d$ за модулем $n$ за схемою Горнера;
    \item \code{window\_exponentiation\_mod\_big\_int} "--- модулярне піднесення
          до степеня віконним методом (4 біти) з таблицею $a^0, \dots, a^{15}$
          (додаткове завдання, варіант 6).
\end{itemize}

Згідно з рекомендаціями завдання, модулярну арифметику реалізовано на базі
редукції Барретта, а піднесення до степеня "--- на базі схеми Горнера.
Усі модулярні функції приймають наперед обчислену структуру \code{BarrettModulus},
тому $\mu$ обчислюється один раз для кожного модуля.
Числа мають фіксовану довжину 2048 біт (32 розряди), а проміжний добуток
$q \cdot \mu$ у редукції Барретта займає до $2k + 2$ розрядів, тому модуль
може мати не більше 15 розрядів (960 біт).

\section{Результати}

Для контролю коректності реалізації алгоритмів (пункт Б завдання) тотожності
перевіряються для декількох випадкових багаторозрядних $a, b, c, n$.
Для генерації випадкових чисел та модулів заданої довжини використано такі
допоміжні функції:

\begin{lstlisting}[style=ruststyle]
fn random_big_int(limbs: usize) -> BigInt {
    let mut rng = rand::rng();
    BigInt {
        limbs: std::array::from_fn(|i| if i < limbs { rng.random() } else { 0 }),
    }
}

fn random_modulus(limbs: usize) -> BarrettModulus {
    let mut n = random_big_int(limbs);
    n.limbs[limbs - 1] |= 1u64 << 63;
    barrett_precompute_big_int(&n).unwrap()
}
\end{lstlisting}

\subsection{Дистрибутивність}

Перевіряється $(a + b) \cdot c \equiv c \cdot (a + b) \equiv a \cdot c + b \cdot c \pmod{n}$
для модулів довжиною від 1 до 15 розрядів.

\begin{lstlisting}[style=ruststyle]
#[test]
fn test_mod_distributivity() {
    for limbs in 1..BIG_INT_LIMBS / 2 {
        let modulus = random_modulus(limbs);
        let a = random_big_int(limbs);
        let b = random_big_int(limbs);
        let c = random_big_int(limbs);

        let sum = add_mod_big_int(&a, &b, &modulus).unwrap();
        let sum_times_c = mult_mod_big_int(&sum, &c, &modulus);
        let c_times_sum = mult_mod_big_int(&c, &sum, &modulus);
        let a_c_plus_b_c = add_mod_big_int(
            &mult_mod_big_int(&a, &c, &modulus).unwrap(),
            &mult_mod_big_int(&b, &c, &modulus).unwrap(),
            &modulus,
        );

        assert_eq!(sum_times_c, c_times_sum);
        assert_eq!(sum_times_c, a_c_plus_b_c);
    }
}
\end{lstlisting}

\subsection{Множення як повторне додавання}

Перевіряється $n \cdot a \equiv \underbrace{a + a + \dots + a}_{n} \pmod{m}$,
де $n$ "--- випадкове число з $[100, 200]$ (не менше за 100).

\begin{lstlisting}[style=ruststyle]
#[test]
fn test_mod_mult_by_repeated_addition() {
    for limbs in 1..BIG_INT_LIMBS / 2 {
        let modulus = random_modulus(limbs);
        let a = random_big_int(limbs);
        let count = rand::rng().random_range(100..=200);

        let mut sum = u64_to_big_int(0);
        for _ in 0..count {
            sum = add_mod_big_int(&sum, &a, &modulus).unwrap();
        }

        assert_eq!(
            mult_mod_big_int(&u64_to_big_int(count), &a, &modulus),
            Some(sum)
        );
    }
}
\end{lstlisting}

\subsection{Теорема Ейлера}

Перевіряється $a^{\varphi(n)} \equiv 1 \pmod{n}$ (за умови $\gcd(a, n) = 1$)
обома методами піднесення до степеня. Для простого $n$ (число Мерсенна
$2^{521} - 1$) $\varphi(n) = n - 1$:

\begin{lstlisting}[style=ruststyle]
#[test]
fn test_euler_theorem_prime() {
    let n = hex_to_big_int(PRIME_521).unwrap();
    let modulus = barrett_precompute_big_int(&n).unwrap();
    let phi = subtract_big_int(&n, &u64_to_big_int(1)).unwrap();

    for _ in 0..3 {
        let a = random_big_int(modulus.k);

        assert_eq!(gcd_big_int(&a, &n), Some(u64_to_big_int(1)));
        assert_eq!(
            horner_exponentiation_mod_big_int(&a, &phi, &modulus),
            Some(u64_to_big_int(1))
        );
        assert_eq!(
            window_exponentiation_mod_big_int(&a, &phi, &modulus),
            Some(u64_to_big_int(1))
        );
    }
}
\end{lstlisting}

Для $n = 3^k$, $k \in \{1, 40, 200, 605\}$, та довільного $a$, що не ділиться на 3,
$\varphi(3^k) = 3^k - 3^{k-1} = 2 \cdot 3^{k-1}$, отже, має виконуватись
$a^{2 \cdot 3^{k-1}} \equiv 1 \pmod{3^k}$:

\begin{lstlisting}[style=ruststyle]
#[test]
fn test_euler_theorem_power_of_three() {
    let three = u64_to_big_int(3);
    for k in [1, 40, 200, 605] {
        let n = horner_exponentiation_big_int(&three, &u64_to_big_int(k));
        let previous_power = horner_exponentiation_big_int(&three, &u64_to_big_int(k - 1));
        let phi = subtract_big_int(&n, &previous_power).unwrap();
        let modulus = barrett_precompute_big_int(&n).unwrap();

        for _ in 0..3 {
            let mut a = random_big_int(modulus.k);
            if divide_mod_big_int(&a, &three).unwrap().1 == u64_to_big_int(0) {
                a = add_big_int(&a, &u64_to_big_int(1));
            }

            assert_eq!(gcd_big_int(&a, &n), Some(u64_to_big_int(1)));
            assert_eq!(
                horner_exponentiation_mod_big_int(&a, &phi, &modulus),
                Some(u64_to_big_int(1))
            );
            assert_eq!(
                window_exponentiation_mod_big_int(&a, &phi, &modulus),
                Some(u64_to_big_int(1))
            );
        }
    }
}
\end{lstlisting}

\subsection{Власні тести}

Окрім тотожностей із завдання, перевіряються:
\begin{itemize}
    \item результати всіх модулярних операцій для фіксованих 950-бітних чисел
          у порівнянні з наперед обчисленими значеннями;
    \item збіг редукції Барретта з остачею від ділення для випадкових $x < \beta^{2k}$;
    \item $(a + b) - b \equiv a \pmod{n}$;
    \item $\gcd(a, b) \cdot \operatorname{lcm}(a, b) = a \cdot b$ для чисел зі спільним дільником;
    \item збіг віконного методу з методом Горнера для показників від 0 до 299;
    \item граничні випадки: нульовий та задовгий модуль, задовгий аргумент
          редукції, НСД і НСК з нулем, нульовий показник степеня.
\end{itemize}

\section{Підтвердження проходження тестів}

Нижче наведено результат запуску \code{cargo test} для тестів цього практикуму;
усі 62 тести бібліотеки пройшли успішно.

\begin{lstlisting}[basicstyle=\ttfamily\scriptsize]
test tests::test_add_mod_big_int ... ok
test tests::test_add_mod_big_int_wrap ... ok
test tests::test_add_subtract_mod_inverse ... ok
test tests::test_barrett_precompute_invalid_modulus ... ok
test tests::test_barrett_reduction ... ok
test tests::test_barrett_reduction_division_consistency ... ok
test tests::test_barrett_reduction_random_division_consistency ... ok
test tests::test_barrett_reduction_too_long ... ok
test tests::test_euler_theorem_power_of_three ... ok
test tests::test_euler_theorem_prime ... ok
test tests::test_gcd_a_greater_or_equal_b ... ok
test tests::test_gcd_a_less_than_b ... ok
test tests::test_gcd_lcm_product ... ok
test tests::test_gcd_zero ... ok
test tests::test_horner_exponentiation_mod_big_int ... ok
test tests::test_horner_exponentiation_mod_big_int_zero_exponent ... ok
test tests::test_kill_big_int_last_limbs ... ok
test tests::test_lcm_big_int ... ok
test tests::test_lcm_zero ... ok
test tests::test_mod_distributivity ... ok
test tests::test_mod_mult_by_repeated_addition ... ok
test tests::test_mult_mod_big_int ... ok
test tests::test_square_mod_big_int ... ok
test tests::test_subtract_mod_big_int_a_greater_than_b ... ok
test tests::test_subtract_mod_big_int_a_less_than_b ... ok
test tests::test_window_exponentiation_mod_big_int ... ok
test tests::test_window_exponentiation_mod_big_int_zero_exponent ... ok
test tests::test_window_horner_exponentiation_mod_consistency ... ok

test result: ok. 62 passed; 0 failed; 0 ignored; 0 measured; 0 filtered out; finished in 1.93s
\end{lstlisting}

\section{Результати бенчмарків}

Середній час виконання реалізованих операцій (пункт В завдання) наведено
в таблиці~\ref{tab:benchmarks} у наносекундах на одну операцію.
Модуль і показники степеня "--- випадкові 960-бітні числа, операнди "--- випадкові
лишки за цим модулем (модуль один для всіх вимірювань, $\mu$ обчислено до початку
вимірювань).
Редукції застосовано до добутку двох операндів.
Для операцій за модулем виконано 1000 прогрівальних запусків і 10000 вимірювань,
для НСД і НСК "--- 100 і 1000, для піднесення до степеня "--- 10 і 100.
Результати отримано командою \code{cargo run -{}-release}.

\begin{table}[htbp]
    \centering
    \begin{tabular}{|l|r|}
        \hline
        \textbf{Операція} & \textbf{Час, нс/операцію} \\
        \hline
        Редукція Барретта & 2769.9 \\
        Редукція діленням з остачею & 33751.4 \\
        Додавання за модулем & 64.4 \\
        Віднімання за модулем & 70.8 \\
        Множення за модулем & 4182.5 \\
        Піднесення до квадрата за модулем & 3289.2 \\
        НСД & 95071.5 \\
        НСК & 131629.7 \\
        Піднесення до степеня за модулем методом Горнера & 5092995.8 \\
        Піднесення до степеня за модулем віконним методом (4 біти) & 4104750.0 \\
        \hline
    \end{tabular}
    \caption{Середній час виконання операцій}
    \label{tab:benchmarks}
\end{table}

Редукція Барретта приблизно в 12 разів швидша
за редукцію діленням з остачею.
Для аргументів, менших за модуль, редукція виконує лише одне порівняння, тому
додавання та віднімання за модулем займають десятки наносекунд, а час множення
та піднесення до квадрата за модулем складається з однієї редукції Барретта
та самої операції.

Віконний метод приблизно на 19\% швидший за метод Горнера.
Для 960-бітного показника обидва методи виконують близько 960 піднесень до квадрата,
але метод Горнера виконує множення для кожного одиничного біта (близько 480),
а віконний метод "--- для кожного ненульового вікна (близько 225, плюс 15 множень
на побудову таблиці).

\end{document}
